\section{持续学习与记忆：上线后的模型不该死掉}

对谈后半段转向研究问题，密度明显变高。谢天宝提出，过去一年 GUI Agent 很多工作更像工程补坑：补数据、补 infra、补 evaluation、补环境。如果要寻找更有范式变化的方向，online learning、continual learning、test-time training 和 self-evolving 可能是一两年内的重要问题。

这个判断的背景是，今天的模型训练范式仍然很“离线”：pre-training、SFT、RL，然后上线推理。上线之后，模型基本不变。即使每天有海量用户交互，模型也不会实时从这些交互中获得能力增益。最多是日志回流、离线清洗、再训练、再部署。问题在于真实产品数据噪声极大，用户反馈稀疏而主观，直接变成训练数据并不容易。

Cursor 的 online RL 被拿出来作为一个信号：当产品场景有高频反馈、明确目标和快速部署 infrastructure 时，模型可以更快进入“产品运营即学习”的循环。但 GUI Agent 比 tab prediction 难得多。它面对的是长程任务、开放动作空间、多工具调用、稀疏 reward 和复杂失败归因。

\begin{importantbox}{上线后的模型不该只是推理工具}
如果模型每天与大量用户交互，却不能从交互中更新知识、偏好、策略或记忆，那么推理过程就只是在消耗算力，而没有产生能力复利。持续学习的目标，是让真实交互成为系统资产，而不是日志垃圾。
\end{importantbox}

Memory 讨论是这一段的核心。谷雨把它拆成三个问题：memory 怎么表示，怎么更新，怎么使用。这个拆法很关键，因为很多产品把 memory 简化成“把历史塞进上下文”或“做一个向量库”，但真正的长期记忆远不止检索。它要决定哪些信息值得长期保留，如何压缩，如何被用户编辑，如何被任务调用，如何避免旧信息污染新决策。

嘉宾们在参数记忆与离散记忆之间来回讨论。参数记忆更新成本高、可塑性差、难以审计；离散记忆如文本、SOP、workflow、knowledge graph 更可编辑、可检查，但也容易信息损耗、结构僵硬、context management 困难。尚晏仪倾向于很多知识要用文本表示，因为文本能进入组织协作；谢天宝则提醒，人脑本身不是纯符号系统，最终可能仍要回到 sub-symbolic 形式。这不是分歧，而是说明未来系统大概率是混合的。

\begin{knowledgebox}{Memory 的三问}
\begin{tabularx}{\textwidth}{>{\raggedright\arraybackslash}p{0.18\textwidth} Y}
\toprule
问题 & 关键含义 \\
\midrule
如何表示 & 参数、文本、向量、SOP、workflow、knowledge graph 各有不同的可塑性、可审计性和调用成本。 \\
如何更新 & 从交互中抽取什么，谁来确认，错误记忆如何删除，偏好变化如何覆盖旧记录。 \\
如何使用 & 是用于检索、规划、个性化、权限控制、错误恢复，还是进入下一轮训练。 \\
\bottomrule
\end{tabularx}
\end{knowledgebox}

Test-time training 和 intrinsic reward 的讨论，则指向另一种可能：模型在推理过程中能否从自身分布、confidence、采样路径或交互结果里获得信号。这里的直觉是，模型在 pre-training 和 post-training 中已经吸收了大量知识，正确答案可能在分布里，只是 decoding 或搜索没有把它取出来。这个方向还远未成熟，但它把“推理”和“学习”的边界重新打开了。

\begin{warningbox}{长上下文不是长期记忆}
把所有历史都塞进 context 会遇到成本、噪声、隐私、检索和一致性问题。长期记忆必须能压缩、能编辑、能遗忘、能审计；否则它不是资产，而是越来越重的认知债务。
\end{warningbox}

\subsection{本章小结}

持续学习讨论把 GUI Agent 从一次性工具推向长期系统。上线后的模型若不能从交互中积累能力，就无法形成真正复利。Memory 不是聊天记录，至少包含表示、更新和使用三件事。未来更可能是参数记忆、离散记忆、workflow 与 test-time adaptation 的混合，而不是单一长上下文或单一向量库。

